\documentclass[10pt,a4paper]{amsart}
\usepackage{amsmath,amssymb,amsthm,graphicx}
% Retained minimal preamble for the Stacks Project chapter more-algebra.tex (More on Algebra).
% The chapter file pinned-original-source.tex carries only \input{preamble} on its first line, so no
% preamble text lives inside the chapter itself.  The Stacks Project uses one single shared preamble for
% every chapter; the snapshot stages it identically as sources/stacks-chow/preamble.tex and
% sources/stacks-etale-cohomology/preamble.tex.  The block below is copied mechanically, line for line,
% from sources/stacks-chow/preamble.tex lines 14-18 (the Xy-pic setup) and 154-217 (the native theorem
% declarations together with the mathematical macros).  Nothing was added, removed or redefined; the
% wrapper class and the amsmath/amssymb/amsthm/graphicx base packages are supplied by the assembler.
% Chapter: more-algebra.tex -- More on Algebra.
% For commutative diagrams we use Xy-pic
\usepackage[all]{xy}

% We use 2cell for 2-commutative diagrams.
\xyoption{2cell}

% Theorem environments.
%
\theoremstyle{plain}
\newtheorem{theorem}[subsection]{Theorem}
\newtheorem{proposition}[subsection]{Proposition}
\newtheorem{lemma}[subsection]{Lemma}

\theoremstyle{definition}
\newtheorem{definition}[subsection]{Definition}
\newtheorem{example}[subsection]{Example}
\newtheorem{exercise}[subsection]{Exercise}
\newtheorem{situation}[subsection]{Situation}

\theoremstyle{remark}
\newtheorem{remark}[subsection]{Remark}
\newtheorem{remarks}[subsection]{Remarks}

\numberwithin{equation}{subsection}

% Macros
%
\def\lim{\mathop{\mathrm{lim}}\nolimits}
\def\colim{\mathop{\mathrm{colim}}\nolimits}
\def\Spec{\mathop{\mathrm{Spec}}}
\def\Hom{\mathop{\mathrm{Hom}}\nolimits}
\def\Ext{\mathop{\mathrm{Ext}}\nolimits}
\def\SheafHom{\mathop{\mathcal{H}\!\mathit{om}}\nolimits}
\def\SheafExt{\mathop{\mathcal{E}\!\mathit{xt}}\nolimits}
\def\Sch{\mathit{Sch}}
\def\Mor{\mathop{\mathrm{Mor}}\nolimits}
\def\Ob{\mathop{\mathrm{Ob}}\nolimits}
\def\Sh{\mathop{\mathit{Sh}}\nolimits}
\def\NL{\mathop{N\!L}\nolimits}
\def\CH{\mathop{\mathrm{CH}}\nolimits}
\def\proetale{{pro\text{-}\acute{e}tale}}
\def\etale{{\acute{e}tale}}
\def\QCoh{\mathit{QCoh}}
\def\Ker{\mathop{\mathrm{Ker}}}
\def\Im{\mathop{\mathrm{Im}}}
\def\Coker{\mathop{\mathrm{Coker}}}
\def\Coim{\mathop{\mathrm{Coim}}}

% Boxtimes
%
\DeclareMathSymbol{\boxtimes}{\mathbin}{AMSa}{"02}

%
% Macros for moduli stacks/spaces
%
\def\QCohstack{\mathcal{QC}\!\mathit{oh}}
\def\Cohstack{\mathcal{C}\!\mathit{oh}}
\def\Spacesstack{\mathcal{S}\!\mathit{paces}}
\def\Quotfunctor{\mathrm{Quot}}
\def\Hilbfunctor{\mathrm{Hilb}}
\def\Curvesstack{\mathcal{C}\!\mathit{urves}}
\def\Polarizedstack{\mathcal{P}\!\mathit{olarized}}
\def\Complexesstack{\mathcal{C}\!\mathit{omplexes}}
% \Pic is the operator that assigns to X its picard group, usage \Pic(X)
% \Picardstack_{X/B} denotes the Picard stack of X over B
% \Picardfunctor_{X/B} denotes the Picard functor of X over B
\def\Pic{\mathop{\mathrm{Pic}}\nolimits}
\def\Picardstack{\mathcal{P}\!\mathit{ic}}
\def\Picardfunctor{\mathrm{Pic}}
\def\Deformationcategory{\mathcal{D}\!\mathit{ef}}

% --- appended 2026-10-12 (Chief): the project own definitions of the reference/slogan/history
% display environments, copied mechanically from the shared Stacks preamble lines 8-12
% (\usepackage{verbatim} plus the three \newenvironment lines, which the chapter preamble keeps
% verbatim from the project preamble). Nothing added, removed or redefined: these environments
% carry no mathematical content and the project itself typesets them as comments. Without these
% three lines a counted statement whose own text uses them fails with Environment slogan undefined.
\usepackage{verbatim}
\newenvironment{reference}{\comment}{\endcomment}
\newenvironment{slogan}{\comment}{\endcomment}
\newenvironment{history}{\comment}{\endcomment}
\usepackage{cleveref}
\crefname{originallemma}{original lemma}{original lemmas}
\Crefname{originallemma}{Original lemma}{Original lemmas}
\makeatletter
\newlabel{lemma-more-relative-pseudo-coherent-is-moot}{{0E1T}{source}{}{sourcepointer-1}{}}
\newlabel{lemma-more-relative-pseudo-coherent-is-moot@cref}{{[originallemma][0][]0E1T}{[1][0][]source}{}{}{}}
\makeatother
\crefname{originallemma}{original lemma}{original lemmas}
\Crefname{originallemma}{Original lemma}{Original lemmas}
\makeatletter
\newlabel{lemma-perfect-over-regular-local-ring}{{09PC}{source}{}{sourcepointer-2}{}}
\newlabel{lemma-perfect-over-regular-local-ring@cref}{{[originallemma][0][]09PC}{[1][0][]source}{}{}{}}
\makeatother
\crefname{originallemma}{original lemma}{original lemmas}
\Crefname{originallemma}{Original lemma}{Original lemmas}
\makeatletter
\newlabel{lemma-check-perfect-stalks}{{068W}{source}{}{sourcepointer-3}{}}
\newlabel{lemma-check-perfect-stalks@cref}{{[originallemma][0][]068W}{[1][0][]source}{}{}{}}
\makeatother
\crefname{originallemma}{original lemma}{original lemmas}
\Crefname{originallemma}{Original lemma}{Original lemmas}
\makeatletter
\newlabel{lemma-perfect-relatively-perfect}{{0DHU}{source}{}{sourcepointer-4}{}}
\newlabel{lemma-perfect-relatively-perfect@cref}{{[originallemma][0][]0DHU}{[1][0][]source}{}{}{}}
\makeatother
\title{Relative perfectness over a flat finitely presented algebra is equivalent to boundedness below plus boundedness of every residue field fibre}
\author{The Stacks Project Authors}
\date{stacks-project:more-algebra.tex:local-snapshot-20261012; GNU FDL 1.2 or later}
\begin{document}
\maketitle
\setcounter{subsection}{1}
\noindent\textit{Editorial note.} All delivered mathematical text is the author's own, copied from the pinned Stacks Project chapter file more-algebra.tex; no statement or proof step was rewritten, reordered or removed. The counted statement is the labeled lemma environment at source lines 22204-22213 (canonical tag 0GHJ) and the counted proof is the proof environment at source lines 22215-22245, which immediately follows it and is its own proof.\par\medskip
\noindent\textit{Editorial note.} The four references are all to earlier results of this chapter, two of which are the local perfectness criterion and the stalk criterion that are themselves the subject of cards of this lane with disjoint coordinates; all four are navigation stubs with their roles recorded in the reasons.\par\medskip
\noindent\textit{Editorial note.} The preamble span runs to line 90, the current end of the retained preamble, which includes the appended environment definitions; the counted text does not use them, and the span is taken whole so that the file version is unambiguous.\par\medskip
\noindent\textit{Editorial note.} The source URL is the source-level Stacks Project URL, as required when more than one card is produced from one source; the per-claim tag is carried by claim\_key and by this note. The counted statement and proof cite no bibliography entry, so this card needs no bibliography.\par\medskip
\noindent\textit{Allowed standard prerequisites (editorial).}\begin{itemize}
\item An R-perfect complex has bounded tor dimension as a complex of R-modules, which gives the easy implication. used at proof lines 22215-22216. Proves one direction of the equivalence.
\item A complex over a flat finitely presented algebra may be viewed as a complex over a polynomial algebra over the base, and pseudo-coherence and relative perfectness are unaffected by this change of algebra because the two have the same image over the base. used at proof lines 22219-22227. Reduces the general case of the statement to the case of a polynomial algebra over the base.
\item For a flat local map whose fibre is regular of finite dimension, a pseudo-coherent complex whose derived fibre at the closed point is bounded below is perfect, and boundedness below of a pseudo-coherent complex is checked on localizations. used at proof lines 22231-22240. Gives perfectness of the localizations at primes of the algebra, whence perfectness of the complex itself.
\end{itemize}\medskip
\ifdefined\hypertarget\hypertarget{sourcepointer-1}{}\fi
\noindent\textit{Source reference.} Original lemma 0E1T: sources/stacks-more-algebra/pinned-original-source.tex, source lines 21885--21885. The reference is invoked at proof lines 22221-22223 to keep pseudo-coherence when the complex is viewed over the polynomial algebra; the reduction to the polynomial case is inside the counted span.\par\medskip
\ifdefined\hypertarget\hypertarget{sourcepointer-2}{}\fi
\noindent\textit{Source reference.} Original lemma 09PC: sources/stacks-more-algebra/pinned-original-source.tex, source lines 20229--20229. The reference is applied at proof lines 22233-22237 to conclude that the localization of the complex at a prime of the algebra is perfect, using the assumption on the fibre; the stalk argument built on it is inside the counted span. The lemma itself is the subject of a separate card of this lane with disjoint coordinates.\par\medskip
\ifdefined\hypertarget\hypertarget{sourcepointer-3}{}\fi
\noindent\textit{Source reference.} Original lemma 068W: sources/stacks-more-algebra/pinned-original-source.tex, source lines 20064--20064. The reference converts perfectness of all localizations, together with boundedness below, into perfectness of the complex, applied at proof lines 22239-22240; the conclusion drawn from it is inside the counted span. The lemma itself is the subject of a separate card of this lane.\par\medskip
\ifdefined\hypertarget\hypertarget{sourcepointer-4}{}\fi
\noindent\textit{Source reference.} Original lemma 0DHU: sources/stacks-more-algebra/pinned-original-source.tex, source lines 21966--21966. The reference is the last step, applied at proof lines 22241-22244 to conclude relative perfectness from perfectness; it is a pointer and the preparation for it is inside the counted span.\par\medskip
\begin{lemma}
\label{lemma-bounded-on-fibres-relatively-perfect}
Let $R \to A$ be a ring map which is flat and of finite presentation.
Let $K \in D(A)$ be pseudo-coherent. The following are equivalent
\begin{enumerate}
\item $K$ is $R$-perfect, and
\item $K$ is bounded below and for every prime ideal $\mathfrak p \subset R$
the object $K \otimes_R^\mathbf{L} \kappa(\mathfrak p)$ is bounded below.
\end{enumerate}
\end{lemma}

\begin{proof}
Observe that (1) implies (2) as an $R$-perfect complex has bounded
tor dimension as a complex of $R$-modules by definition. Let us prove
the other implication.

\medskip\noindent
Write $A = R[x_1, \ldots, x_d]/I$. Denote $L$ in $D(R[x_1, \ldots, x_d])$
the restriction of $K$. By
Lemma \ref{lemma-more-relative-pseudo-coherent-is-moot}
we see that $L$ is pseudo-coherent.
Since $L$ and $K$ have the same image in $D(R)$ we see that
$L$ is $R$-perfect if and only if $K$ is $R$-perfect.
Also $L \otimes_R^\mathbf{L} \kappa(\mathfrak p)$
and $K \otimes_R^\mathbf{L} \kappa(\mathfrak p)$ are the same objects of
$D(\kappa(\mathfrak p))$. This reduces us to the case
$A = R[x_1, \ldots, x_d]$.

\medskip\noindent
Say $A = R[x_1, \ldots, x_d]$ and $K$ satisfies (2).
Let $\mathfrak q \subset A$ be a prime lying over a
prime $\mathfrak p \subset R$. By
Lemma \ref{lemma-perfect-over-regular-local-ring} applied
to $R_\mathfrak p \to A_\mathfrak q$ and the complex $K_\mathfrak q$
using our assumption, we find that $K_\mathfrak q$ is perfect
in $D(A_\mathfrak q)$. Since $K$ is bounded below, we see that
$K$ is perfect in $D(A)$ by 
Lemma \ref{lemma-check-perfect-stalks}.
This implies that $K$ is $R$-perfect by
Lemma \ref{lemma-perfect-relatively-perfect}
and the proof is complete.
\end{proof}

\end{document}
